\section{Анализ подходов к разработке интерфейса интеллектуальной системы}
\label{sec:analysis}

\subsection{Характеристика предметной области}

% Требования:
% - описать предметную область, в которой применяется интеллектуальная система;
% - указать типы решаемых задач и ожидаемые результаты работы системы;
% - зафиксировать ограничения и особенности (объёмы данных, требования по времени ответа и т.п.).

Предметной областью разрабатываемой системы является человеко-машинное взаимодействие в рамках сложных виртуальных сред и имитационных моделей. В частности, проект фокусируется на проблематике создания интеллектуальных пользовательских интерфейсов, основанных на обработке естественного языка и генеративном искусственном интеллекте\cite{Nielsen1993}.

Современная индустрия интерактивных развлечений и виртуальных симуляций демонстрирует устойчивую тенденцию к усложнению игровых механик и увеличению объемов информации, поступающей пользователю в режиме реального времени. Традиционные подходы к организации интерфейсов, базирующиеся на графических элементах, меню и статических текстовых подсказках, часто оказываются недостаточно эффективными для обеспечения иммерсивного опыта. Пользователь вынужден переключать внимание между игровым процессом и интерфейсными элементами, что повышает когнитивную нагрузку и нарушает состояние потока\cite{Csikszentmihalyi1990,Sweller1988}.

В качестве полигона для исследования и разработки выбрана среда Minecraft — песочница с открытым миром, которая характеризуется высокой степенью энтропии событий. В отличие от линейных приложений, действия пользователя здесь не предопределены скриптами, что создает поток неструктурированных данных, который проблематично интерпретировать классическими алгоритмическими методами. Большие языковые модели способны выступать в роли когнитивного модуля, который не просто реагирует на атомарные триггеры, но и формирует целостную картину происходящего, обладает памятью и способен к ролевому моделированию\cite{Brown2020,Park2023}.

Разрабатываемая интеллектуальная система ассистирования предназначена для решения комплекса задач, направленных на улучшение пользовательского опыта\cite{Wang2023}:
\begin{itemize}
    \item агрегация и интерпретация событийного потока, подразумевающая перехват низкоуровневых данных от клиента приложения и их преобразование в семантически значимые описания ситуаций;
    \item контекстно-зависимая генерация речи, требующая реализации механизма принятия решений о целесообразности речевого вмешательства вместо стандартной схемы запрос-ответ;
    \item эмоциональная адаптация, при которой ассистент поддерживает заданный характер, изменяя тональность ответов и стиль речи в зависимости от успеха или неудач пользователя;
    \item мультимодальное представление информации через синхронизированный вывод данных посредством аудиоканала и визуального оверлея.
\end{itemize}

Внедрение интеллектуального ассистента в виртуальную среду должно привести к следующим результатам:
\begin{itemize}
    \item снижение когнитивной нагрузки, так как пользователь получает необходимую информацию через естественную речь, не отвлекаясь на чтение текстовых логов;
    \item повышение иммерсивности за счет замены механических интерфейсных уведомлений на диалог с персонифицированным персонажем;
    \item обеспечение доступности через дублирование информации по разным каналам восприятия.
\end{itemize}

Проектирование систем на базе языковых моделей в контексте реального времени накладывает ряд жестких технических и логических ограничений\cite{Xi2023}:
\begin{itemize}
    \item требования по времени отклика, согласно которым задержка между игровым событием и реакцией ассистента не должна превышать 1.5–2 секунды для сохранения иллюзии живого общения;
    \item ограничение контекстного окна языковой модели, требующее применения механизмов буферизации и суммаризации контекста для выделения только релевантных фактов из длительной сессии;
    \item склонность моделей к генерации правдоподобной, но ложной информации (галлюцинациям), что минимизируется через системные промпты и валидацию данных;
    \item большие объемы неструктурированных данных, генерируемых игровым клиентом, которые делают экономически и технически нецелесообразной прямую передачу всех событий в модель без предварительной фильтрации;
    \item этические аспекты, связанные с необходимостью фильтрации контента при имитации антагонистических характеров.
\end{itemize}


\subsection{Участники взаимодействия и их задачи}

% Требования:
% - перечислить участников: конечные пользователи, внешние системы, устройства, внутренние подсистемы;
% - описать роли и цели каждого участника (какие задачи решает, какие данные вводит/получает);
% - определить контекст использования интерфейса для каждой роли (рабочее место, устройства, сценарии).


Архитектура разрабатываемого программно-аппаратного комплекса предполагает взаимодействие между человеком-оператором, локальными программными компонентами и удаленными облачными сервисами. Участников информационного обмена можно классифицировать на четыре ключевые группы: конечный пользователь, внешняя виртуальная среда, внутренние подсистемы обработки данных и внешние когнитивные сервисы.

\textbf{Конечный пользователь (оператор системы)}

Основным выгодоприобретателем является пользователь, осуществляющий взаимодействие с виртуальной средой. Его роль заключается в генерации событийного потока через игровые механики и потреблении аудиовизуальной обратной связи от ассистента. Задачи пользователя в системе:
\begin{itemize}
    \item осуществление игровой деятельности, инициирующей поток событий (перемещение, взаимодействие с объектами, бой, строительство);
    \item настройка параметров личности ассистента и системных промптов через веб-интерфейс для персонализации опыта;
    \item восприятие синтезированной речи и визуальных индикаторов оверлея для корректировки своих действий в виртуальной среде.
\end{itemize}
Контекст использования предполагает наличие стационарного персонального компьютера, оснащенного устройствами вывода звука (наушники или акустическая система) и, опционально, вторым монитором для отображения панели управления.

\textbf{Клиентское приложение виртуальной среды (Minecraft)}

Внешняя система, выступающая генератором первичных данных. Роль данного участника — выполнение симуляции виртуального мира и фиксация изменений состояния. Задачи компонента:
\begin{itemize}
    \item перехват внутриигровых триггеров посредством модификации кода игры;
    \item сериализация событий в структурированный формат (JSON) с добавлением метаданных (временная метка, координаты, тип события);
    \item передача данных на сервер обработки в режиме реального времени;
    \item исполнение команд обратной связи, поступающих от интеллектуальной системы (например, изменение погоды или выдача предметов).
\end{itemize}

\textbf{Серверная часть системы (ядро обработки)}

Центральный узел, оркестрирующий взаимодействие всех компонентов. Его цель — преобразование низкоуровневых сигналов игры в семантически значимый контент. Задачи подсистемы:
\begin{itemize}
    \item агрегация и буферизация входящих событий для предотвращения избыточной частоты запросов;
    \item управление контекстом диалога и краткосрочной памятью ассистента;
    \item принятие решений о необходимости реакции на основе анализа текущей ситуации;
    \item маршрутизация запросов к внешним API и передача готового аудиовизуального контента на клиентские интерфейсы (оверлей и веб-панель).
\end{itemize}

\textbf{Внешние сервисы искусственного интеллекта}

Сторонние облачные или локальные решения, обеспечивающие когнитивные функции системы. К ним относятся большие языковые модели (LLM) и сервисы синтеза речи (TTS). Задачи сервисов:
\begin{itemize}
    \item генерация естественного текста ответа на основе переданного системного промпта и контекста событий;
    \item классификация намерений пользователя и эмоциональная окраска реплик;
    \item преобразование текстовых строк в аудиопоток с заданными характеристиками тембра и скорости речи.
\end{itemize}

Взаимодействие между участниками осуществляется асинхронно, что позволяет поддерживать высокую отзывчивость интерфейса даже при задержках в обработке данных на стороне внешних сервисов.

\subsection{Анализ существующих систем и интерфейсов}
Проектирование нового интеллектуального интерфейса требует детального рассмотрения существующих решений в области автоматизации взаимодействия с виртуальными средами. На текущем этапе развития технологий можно выделить несколько классов систем, частично реализующих функционал ассистирования: императивные навигационные боты, исследовательские автономные агенты на базе больших языковых моделей и встроенные системы доступности.

\textbf{Обзор аналогичных систем}

% Требования:
% - выбрать не менее 2–3 систем/продуктов схожей тематики;
% - для каждой системы кратко описать назначение и тип интерфейса (человек–система, система–система, система–устройство).
В качестве объектов анализа выбраны три репрезентативные системы, демонстрирующие различные подходы к организации интерфейса человек-компьютер в контексте виртуальной среды Minecraft:

1. \textbf{Baritone (Pathfinding System).} Это наиболее распространенная система автоматизированной навигации и выполнения рутинных операций. Назначение системы заключается в алгоритмическом расчете оптимальных маршрутов и автоматическом выполнении действий (добыча ресурсов, строительство) без прямого управления персонажем. Тип интерфейса можно охарактеризовать как командно-текстовый (CLI over GUI), где взаимодействие происходит посредством ввода специальных директив в стандартный чат игры\cite{Baritone}.

2. \textbf{Voyager (Embodied Agent).} Экспериментальная система, разработанная исследователями из NVIDIA и Caltech. Это первый воплощенный агент на базе LLM (GPT-4), способный к непрерывному обучению без участия человека. Назначение системы — автономное исследование мира, открытие новых технологий и накопление библиотеки навыков в виде исполняемого кода. Интерфейс системы относится к типу система-система, так как управление осуществляется программно, а роль человека сводится к наблюдению за логами и метриками через консольный вывод, без возможности прямого диалогового вмешательства в реальном времени\cite{Wang2023}.

3. \textbf{Стандартный экранный диктор (Minecraft Narrator).} Встроенный инструмент обеспечения доступности, преобразующий текстовые сообщения чата в речь. Назначение системы — обеспечение инклюзивности для пользователей с нарушениями зрения. Тип интерфейса — односторонний аудио-вывод, жестко привязанный к системным событиям графического интерфейса\cite{MinecraftAccessibility}.

\textbf{Анализ интерфейсов существующих решений}

% Требования:
% - для каждого аналога рассмотреть ключевые сценарии взаимодействия, структуру интерфейса, типовые операции;
% - выделить сильные и слабые стороны (удобство, прозрачность, гибкость, ошибки проектирования интерфейса);
% - использовать корректные ссылки на источники (по требованиям ВАК).
Детальный анализ пользовательского взаимодействия (UX) с перечисленными системами позволяет выявить фундаментальные ограничения текущих подходов и сформулировать требования к проектируемому мультимодальному ассистенту.

\textbf{Анализ интерфейса системы Baritone}

Ключевой сценарий взаимодействия с Baritone строится на использовании префиксов (обычно символа \#) в игровом чате. Пользователь должен помнить синтаксис команд, например, директива \texttt{\#mine diamond\_ore} инициирует поиск алмазов.
Структура интерфейса представляет собой гибрид графического отображения игрового мира и текстовой командной строки. Визуализация намерений бота реализуется через рендеринг цветных линий маршрута поверх изображения.

Сильные стороны решения заключаются в высокой детерминированности и точности исполнения: пользователь получает полный контроль над действиями аватара. Оверлей с траекторией движения обеспечивает прозрачность алгоритмических решений системы.

Однако проявляются критические ошибки проектирования интерфейса для массового пользователя:
\begin{itemize}
    \item высокий порог вхождения, обусловленный необходимостью запоминания сотен команд и их параметров, что создает чрезмерную когнитивную нагрузку;
    \item отсутствие контекстной гибкости, так как система не понимает абстрактных запросов на естественном языке (например, «найди безопасное место») и требует четкой императивной формулировки;
    \item модальность ввода блокирует управление персонажем, поскольку для ввода команды необходимо открыть окно чата, что делает невозможным оперативное реагирование в критических ситуациях (например, во время боя).
\end{itemize}

\textbf{Анализ интерфейса агента Voyager}

Взаимодействие с системой Voyager происходит преимущественно в режиме наблюдателя. Пользователь запускает скрипт инициализации, после чего система работает автономно, генерируя код навыков на языке JavaScript. Структура интерфейса для оператора представлена набором текстовых логов в терминале и визуальным наблюдением за действиями бота в окне игры.

Сильной стороной является полная автоматизация когнитивных процессов: система сама ставит цели и достигает их. Это демонстрирует потенциал использования LLM для сложного планирования.

С точки зрения пользовательского интерфейса ассистента, данное решение имеет существенные недостатки:
\begin{itemize}
    \item отсутствие интерактивности, так как архитектура Voyager не предусматривает канала обратной связи от пользователя к агенту в реальном времени;
    \item непрозрачность принятия решений для неквалифицированного пользователя, который видит лишь конечный результат действий, но не понимает мотивацию агента без анализа программного кода;
    \item задержка реакции, обусловленная итеративным процессом генерации и отладки кода, что неприемлемо для задач оперативного ассистирования.
\end{itemize}

\textbf{Анализ интерфейса экранного диктора (Narrator)}

Сценарий использования сводится к пассивному прослушиванию синтезированной речи, дублирующей текстовые сообщения других игроков и системные уведомления. Интерфейс не имеет визуального представления, управление осуществляется через меню настроек.

Преимуществом является простота интеграции и низкое потребление ресурсов. Однако функциональность системы ограничена прямым чтением текста без семантического анализа.

Слабые стороны включают:
\begin{itemize}
    \item отсутствие фильтрации контекста, приводящее к зашумлению аудиоканала (информационный спам), когда диктор озвучивает технические логи или повторяющиеся сообщения;
    \item монотонность и отсутствие эмоциональной окраски, что снижает вовлеченность и не позволяет использовать систему как полноценного компаньона;
    \item отсутствие мультимодальности, так как аудиопоток не подкрепляется визуальными образами или анимацией, что затрудняет восприятие информации в шумной обстановке.
\end{itemize}




\subsection{Анализ требований к интерфейсу}

На основе исследования предметной области и ролевой модели участников взаимодействия определены требования к интерфейсу интеллектуальной системы ассистирования.

\textbf{Функциональные требования}

Функциональность интерфейса распределена между конечным пользователем (игроком) и администратором системы:
\begin{itemize}
    \item {интерфейс игрового клиента (мод)} должен обеспечивать автоматический сбор и передачу структурированных данных о состоянии аватара (HP, уровень сытости) и окружающей среды (наличие мобов, время суток) на сервер.
    \item {интерфейс обратной связи (внутри игры)} должен обеспечивать визуальное отображение текстовых советов и предупреждений, полученных от сервера, в удобной для восприятия форме (чат или графический оверлей).
    \item {веб-интерфейс администратора} должен обеспечивать визуализацию потока игровых событий в режиме реального времени посредством протокола WebSocket для оперативного мониторинга состояния системы.
    \item {интерфейс управления конфигурацией} должен обеспечивать возможность изменения рабочих параметров сервера (например, \textit{analysis\_interval}) через веб-форму с немедленной синхронизацией настроек с игровым клиентом.
    \item {система должен обеспечивать} автоматическую валидацию входящих сообщений на соответствие схеме данных \textit{LogEntry} перед их отображением или анализом.
\end{itemize}

\textbf{Нефункциональные требования}

Для обеспечения стабильной работы в условиях высокодинамичной среды Minecraft к системе предъявляются следующие требования:
\begin{itemize}
    \item надежность: для предотвращения сбоев при передаче данных необходимо использовать строгую валидацию типов через \textit{Pydantic}. Не допускается потеря событий при кратковременных разрывах соединения (необходима буферизация на стороне клиента).
    \item удобство и доступность: интерфейс админ-панели следует реализовывать на базе стандартных веб-технологий (HTML5/JS), обеспечивающих доступность через любой современный браузер без установки дополнительного ПО.
    \item ограничения среды: клиентская часть должна быть совместима со средой \textit{Fabric 1.20.1} и \textit{Java 17}\cite{FabricAPI}. Не допускается блокировка основного игрового потока (Render Thread) при выполнении сетевых запросов — все операции обмена данными должны быть асинхронными.
    \item расширяемость: архитектуру системы следует проектировать таким образом, чтобы замена \textit{rule-based} логики на \textit{LLM-модуль} не требовала переработки интерфейсов взаимодействия.
\end{itemize}

\textbf{Требования к «интеллектуальности» интерфейса}

Интеллектуальная составляющая интерфейса в рамках курсового проекта реализуется через событийный анализ и адаптивную подачу информации:
\begin{itemize}
    \item адаптивная фильтрация (Троттлинг): интерфейс должен реализовывать логику подавления избыточных уведомлений. Не допускается повторная выдача одного и того же совета (например, о наступлении ночи) чаще, чем установлено заданным лимитом времени (от 10 до 60 секунд в зависимости от типа события).
    \item приоритизация контекста: при возникновении нескольких триггеров одновременно (например, голод и низкое здоровье) система должна выбирать приоритетный совет. Логика выбора должна базироваться на степени угрозы для выживания игрового персонажа.
    \item интеллектуальная обработка состояний: система не просто дублирует игровые логи, а интерпретирует их. Например, на основе типа события \textit{near\_hostile} и данных о расстоянии интерфейс должен формировать предиктивное предупреждение об угрозе атаки.
    \item уровень сложности: для текущего прототипа устанавливается уровень \textit{экспертной системы на базе правил (rule-based)}. Интеллектуальность подтверждается наличием четких алгоритмов обработки контекста, описанных в коде сервера (\textit{advice.py}), а не случайной выдачей фраз.
\end{itemize}

\subsection{Вывод по разделу анализа}

На основе проведенного анализа можно сделать следующие выводы:
\begin{itemize}
    \item предметная область характеризуется высокой энтропией данных, что требует реализации посредника между сырыми событиями игры и пользователем для снижения когнитивной нагрузки.
    \item анализ существующих решений показал, что большинство аналогов либо излишне сложны (Baritone), либо не обеспечивают интерактивного взаимодействия в реальном времени (Voyager), что подтверждает актуальность разработки легковесного ассистента.
    \item в качестве архитектурной основы выбрана клиент-серверная модель: клиент на базе \textit{Fabric} для сбора данных и сервер на \textit{FastAPI} для анализа контекста и управления настройками.
    \item принято решение использовать \textit{rule-based} подход для генерации советов на этапе MVP, что обеспечит детерминированность поведения системы и заложит фундамент для последующей интеграции языковых моделей.
\end{itemize}